\section{Introduction}
\label{sec:intro}

A \emph{hammer} is a proof-automation tool that, given a goal stated in the
logic of a proof assistant and the ambient library of previously established
facts, selects a set of plausibly relevant premises, translates the goal and
the premises into the input logic of external automated theorem provers
(ATPs), runs the ATPs, and finally reconstructs any proof found inside the
proof assistant~\cite{BlanchetteEtAl2016qed}.  For proof assistants based on
dependent type theory, the reference tool is
CoqHammer~\cite{CzajkaKaliszyk2018,CoqHammer} for Rocq (formerly Coq), whose
translation maps the Calculus of Inductive Constructions (CIC) into untyped
first-order logic (FOL).

The efficiency of the whole enterprise hinges on the translation being
\emph{shallow}~\cite{Czajka2018}: a source proposition must become a
first-order \emph{formula}, with the source connectives and quantifiers
mapped directly to their first-order counterparts, so that the ATP's highly
optimised handling of first-order primitives --- clausification,
resolution, superposition, term indexing --- does real logical work on the
translated problem.  The alternative, a \emph{deep} embedding in which
propositions become first-order \emph{terms} related by an inhabitation
predicate $\mathsf{T}(x,\tau)$, forces the ATP to simulate the object logic
one bridge-axiom instantiation at a time and to synthesise first-order terms
standing for proof terms; this is well understood to be dramatically less
effective (see the discussion in \cite[\S 1.1]{Czajka2018} and
\cref{sec:shallow} below).

For the \emph{propositional} skeleton of a statement, shallowness is
achieved by the translation of \cite{CzajkaKaliszyk2018}: universally
quantified implications become guarded first-order formulas, equality
becomes first-order equality, and so on.  Two related parts of the
translation, however, have remained essentially deep, and both concern
exactly the places where dependent type theory diverges from first-order
logic:

\begin{enumerate}[leftmargin=2em]
\item \emph{Definitions by pattern matching.}  The translation of
  \cite{CzajkaKaliszyk2018} renders a match
  $\tcase{x}{\cdot}{\csO \Rightarrow t_1 \mid \csS\,y \Rightarrow t_2}$,
  lifted to a fresh constant $F$, as a single \emph{guarded disjunctive
  axiom}
  \[
    \forall x.\ \nat(x) \to
      \bigl(x \feq \csO \wedge F(x) \feq t_1'\bigr) \vee
      \bigl(\exists y.\ \nat(y) \wedge x \feq \csS(y) \wedge F(x) \feq t_2'\bigr),
  \]
  a formula that clausifies into many long non-Horn clauses with Skolem
  functions and that superposition provers cannot orient into rewrite
  rules.
\item \emph{Types with propositional content.}  A constant whose type
  mixes computation with specification --- the running example, from the
  CoqHammer problem list, is
  \begin{align*}
    &h : \forall x\,y\,z{:}\nat.\; x = y \wedge y = z \to
      \{u : \nat \mid x = u\}\\
    &h\;x\;y\;z\;p \defeq
      \mathsf{match}\ p\ \mathsf{with}\ \mathsf{conj}\ p_1\ p_2
      \Rightarrow{}\\
    &\qquad\qquad
      \mathsf{exist}\ (\lambda u.\, x = u)\ z\ (\mathsf{eq\_trans}\ p_1\ p_2)
  \end{align*}
  --- receives \emph{no} usable axiom at all: the match on the proof $p$ is
  translated to an opaque constant, and the type is translated to an opaque
  atom $\mathsf{HasType}(h, \mathsf{sig}(\nat, P))$ that no other axiom
  connects to anything.  The information is present but buried inside an
  inhabitation atom: a local relapse into the deep embedding.
\end{enumerate}

Both phenomena have a common diagnosis, already conjectured in
\cite[\S 9]{CzajkaKaliszyk2018}: the translation should be \emph{factored
through the intermediate representation of program extraction} in the sense
of Letouzey~\cite{Letouzey2004,Letouzey2003,Letouzey2008}.  This paper
carries the factorisation out precisely and proves it sound.

\subsection{The architecture}
\label{sec:intro-arch}

The slogan is:
\begin{quote}
  \emph{translation = program extraction + specification extraction.}
\end{quote}
For a definition $c \defeq t : \tau$ in the source environment, the
translation emits two groups of first-order axioms.

\paragraph{Program extraction.}  The body $t$ is erased, following
Letouzey's function $\erasef$, into a term of an untyped functional
calculus $\lbox$ (\cref{sec:lambdabox}): proofs and propositional arguments
disappear, matches on propositional singletons (conjunctions, equations,
$\bot$) collapse to their unique branch, and refinement values
$\mathsf{exist}\;a\;p$ unbox to $a$.  The erased program is then compiled
--- by lambda-lifting, case flattening and pattern-match compilation
(\cref{sec:translation}) --- into a set of \emph{defining equations}, one
per constructor pattern.  For the standard addition function the output is
literally the recursive program, read as a first-order rewrite system:
\[
  \forall m.\ \Ihat{\mathsf{add}}(\csO, m) \feq m,
  \qquad
  \forall y\,m.\ \Ihat{\mathsf{add}}(\csS(y), m) \feq \csS(\Ihat{\mathsf{add}}(y, m)).
\]
These are orientable unit equations --- demodulators for superposition
provers, ideal trigger material for SMT solvers --- in place of the guarded
disjunction above.  For $h$, erasure yields $\erase{t} = z$ and hence the
single equation $\forall x\,y\,z.\ \Ihat{h}(x,y,z) \feq z$.

\paragraph{Specification extraction.}  The type $\tau$ is translated by a
realizability reading $\guardf$ in the tradition of
Kleene~\cite{Kleene1945} and Paulin-Mohring~\cite{PaulinMohring1989},
defined structurally on types: informative quantification becomes guarded
first-order quantification, propositional premises become implications with
the (erased) proof argument dropped, and --- the essential new clause ---
a refinement type $\refty{x{:}\sigma}{\varphi}$ is expanded \emph{inline,
at every occurrence}, into the conjunction of the carrier guard and the
payload formula.  For $h$ this produces the specification axiom
\begin{multline*}
  \forall x\,y\,z.\ \nat(x) \to \nat(y) \to \nat(z) \to{}\\
    \bigl(x \feq y \wedge y \feq z \;\to\;
      \nat(\Ihat{h}(x,y,z)) \wedge x \feq \Ihat{h}(x,y,z)\bigr),
\end{multline*}
which is exactly the statement of $h$'s type, as a first-order formula.  No
$\mathsf{HasType}$ atom, no $\mathsf{sig}$, no $\mathsf{exist}$ and no
projection appears anywhere in the output: the specification is shallow in
precisely the sense of \cite[\S 1.1]{Czajka2018}.

The two components are the two halves of extraction correctness: the
equations describe the extracted program; the specification axiom is the
statement that the program \emph{realizes} its type, which is Letouzey's
simulation predicate $\sem{\tau}$ \cite[\S 2.4]{Letouzey2004} applied to
the extracted constant.

\subsection{Contributions}
\label{sec:intro-contrib}

This paper makes the architecture precise and proves it correct.
Specifically:

\begin{enumerate}[leftmargin=2em]
\item We define a source calculus $\CICm$ (\cref{sec:calculus}): a
  monomorphic, first-order fragment of CIC with inductive datatypes,
  primitive refinement types, inductive predicates, propositional
  singleton eliminations into informative positions (conjunction,
  equality, absurdity), structural fixpoints, and full first-order
  intuitionistic logic with induction at the proof level.  The calculus is
  \emph{stratified}: terms, proofs, types, propositions and predicates are
  separate syntactic categories, so the erasure zones of extraction are
  syntactic rather than derived from sort analysis.  \Cref{sec:calculus-cic}
  explains how the fragment arises from full CIC by the
  instance-classification analysis implemented in CoqHammer.
\item We define the erasure $\erasef$ from $\CICm$ into an untyped calculus
  $\lbox$ with constructors, case and fixpoints (\cref{sec:lambdabox}), and
  prove the two syntactic pillars of program extraction in this setting:
  erasure commutes with substitution, and erasure maps source conversion to
  $\lbox$-conversion.  $\lbox$ is an orthogonal rewrite system; we prove its
  confluence in \cref{sec:confluence} by the Tait--Martin-L{\"o}f/Takahashi
  method.
\item We define the translation (\cref{sec:translation}): the compilation of
  erased programs to per-constructor pattern equations with case flattening
  and lambda-lifting; the specification extraction $\guardf$, the
  proposition translation $\Ff$ and the term translation $\compf$; and the
  emitted theory $\ThE$ for an environment $\Env$ and goal $G$.  All three
  translation functions are structural --- there is no normalisation
  requirement on the source --- and produce formulas built exclusively from
  first-order connectives, per-datatype guard predicates, per-instance
  inductive-predicate atoms, and equality.
\item We prove soundness (\cref{sec:semantics,sec:soundness}).  From any
  well-formed definitional environment we construct a first-order structure
  $\ML$ whose domain is the set of closed $\lbox$ terms modulo conversion
  --- crucially including ``junk'' elements denoting no source value ---
  and a proof-irrelevant classical truth semantics $\pv{\cdot}$ for source
  propositions, and show: every axiom of $\ThE$ holds in $\ML$
  (\cref{thm:env-validity}), and $\ML \vDash \F{G}$ implies $\pv{G} = 1$
  (\cref{lem:coincidence}).  Hence if $\ThE \provesfol \F{G}$ then $G$ is
  true in the semantics; in particular no goal refutable in $\CICm$ becomes
  provable after translation, and $\ThE$ is consistent
  (\cref{thm:soundness,cor:consistency}).  The technical core is a
  fundamental lemma (\cref{lem:fundamental}) proved by induction on typing
  derivations, with a rank-indexed treatment of structural fixpoints in the
  style of type-based termination~\cite{BartheEtAl2004,Gimenez1995}.
\item We settle, rigorously, the design question that motivated the
  factorisation (\cref{sec:casesplit}): the guarded disjunctive case axiom
  of \cite[\S 5.2]{CzajkaKaliszyk2018} is \emph{false} in $\ML$ if its
  guard is omitted --- we exhibit the junk countermodel --- while the
  per-constructor equations hold at all elements of $\ML$, junk included,
  so they need no guards at all; and the two axiomatisations are
  interderivable in FOL relative to the structural axioms (inversion,
  discrimination, injectivity) that the translation emits anyway.  The
  choice between them is therefore purely operational, and
  \cref{sec:casesplit-cnf} quantifies the operational difference in clause
  shapes.
\item We articulate the shallowness of the translation as a checklist of
  verifiable properties and verify them clause by clause
  (\cref{sec:shallow}), and we delimit precisely what is \emph{not} covered
  --- polymorphism, indexed inductive families, well-founded recursion,
  large elimination --- together with the known hard constraint for
  extending the architecture to well-founded recursion
  (\cref{sec:limitations}).
\end{enumerate}

\subsection{What soundness means here}
\label{sec:intro-soundness}

Following automated-reasoning terminology~\cite{Czajka2018,
BlanchetteEtAl2016}, a hammer translation is \emph{sound} when derivability
of the translated goal from the emitted theory guarantees that the source
goal holds.  Concretely, in the form we prove (\cref{thm:soundness}): if
$\ThE \provesfol \F{G}$ then $G$ is true in a fixed semantics of the source
development (\cref{sec:semantics}).  This is weaker than exhibiting a source
\emph{proof} of $G$ --- a distinction we return to below --- but it already
excludes the failure that matters: no goal that is false in that semantics,
and in particular none refutable in the source ($\Env \vdash \neg G$), has a
derivable translation (\cref{cor:no-refutable}).  (The converse property,
that provable
goals remain provable, is \emph{completeness}; hammers deliberately trade
it away, and the present translation is complete only on a fragment.  See
\cref{sec:limitations-completeness}.)  Since the final arbiter in a hammer
is proof reconstruction inside the proof assistant, soundness is not a
safety requirement but a quality requirement --- an unsound translation
wastes ATP time on phantom proofs and poisons premise-selection learning
\cite{BlanchetteEtAl2016qed} --- and, more importantly for this paper, the
soundness \emph{proof} is the specification of the translation: it
identifies exactly which axioms may be emitted, which guards are mandatory,
and why.

We prove soundness model-theoretically, against a semantics that interprets
source propositions classically and proof-irrelevantly
(\cref{sec:semantics}).  Both choices are forced: the ATPs are classical,
and the translation erases proofs, so it validates proof irrelevance by
construction --- the corresponding necessity was isolated for the core
embedding in \cite[Example 52]{Czajka2018}.  The semantics is nevertheless
\emph{standard} on the first-order data of interest: the interpretation of
$\nat$ in our model is (isomorphic to) the natural numbers
(\cref{prop:standard}), so for concrete arithmetic-style goals ``true in
the semantics'' means true.

A proof-theoretic soundness theorem in the style of \cite{Czajka2018} ---
reconstructing a source proof term from a first-order derivation --- is
strictly stronger and correspondingly harder in the presence of inductive
types; we leave it as future work (\cref{sec:limitations}), noting that the
model-theoretic statement is the one with operational consequences for a
hammer.

\subsection{Relation to the implemented translation}
\label{sec:intro-impl}

The calculus $\CICm$ and the translation of this paper are an idealisation
of a design for CoqHammer, not a description of its current code: the
current tool emits the guarded disjunctive case axioms of
\cref{sec:casesplit} and translates the type of $h$ to an opaque
inhabitation atom.  The idealisation is faithful in the other direction:
every axiom family emitted by the translation of this paper has a direct
implementation site in the existing translation pipeline, and the fragment
$\CICm$ is the image of CIC under analyses (proposition detection,
singleton classification, refinement classification, monomorphisation)
that are either implemented or specified for implementation.  We indicate
the correspondences where relevant but keep the development self-contained.
